% !TeX root = ../main.tex
\subsection{群作用}\label{sec:group-actions}

% 来源目录：../../NotePhotos/BaiduSyncdisk/algebra/2.4/（相对于本科目目录）。
% 共 4 张照片，完整文件名分别记在对应正文前；文件名编号 55--58 不是原稿页码。
% 覆盖原稿 2.4.1--2.4.15，保留黑笔主体、紫笔说明与原稿的证明／提示范围。

% 来源：微信图片_20261118134403_55_26.jpg，节首、2.4.1--2.4.3 及紫笔说明。
\begin{purplenote}
本节研究的技术将用于后续对非阿贝尔有限群结构的研究。
\end{purplenote}

\begin{definition}[群作用]\label{def:group-action}
设 $G$ 是群，$S$ 是集合。$G$ 在 $S$ 上的一个\textbf{作用}是一个映射
\[
  G\times S\longrightarrow S,\qquad (g,x)\longmapsto gx,
\]
满足对任意 $x\in S$ 及 $g_1,g_2\in G$，都有
\[
  ex=x,\qquad (g_1g_2)x=g_1(g_2x).
\]
给定这样的映射后，称 $G$ 作用于集合 $S$。
\end{definition}

\begin{purplenote}
本节采用左作用，应注意作用与平移的方向；对于一般子群的左陪集，不能直接用右平移定义作用。

对给定的群 $G$ 和集合 $S$，可以有许多不同的作用映射，因此使用 $gx$ 等记号时需要注意上下文。不过，后面研究结构时，许多结论并不依赖作用映射的具体形式。
\end{purplenote}

\begin{example}\label{ex:group-actions-permutations-translations-conjugation}
\begin{enumerate}[label=(\arabic*)]
  \item 对称群 $S_n$ 在集合 $I_n=\{1,\ldots,n\}$ 上有自然作用
  \[
    (\sigma,x)\longmapsto\sigma(x).
  \]

  \item 设 $H\leq G$。由 $G$ 中的乘法定义 $H$ 在集合 $G$ 上的作用
  \[
    (h,x)\longmapsto hx.
  \]
  对固定的 $h\in H$，映射 $x\mapsto hx$ 称为由 $h$ 给出的\textbf{左平移}（left translation）。若 $K\leq G$，$S$ 是 $K$ 在 $G$ 中的所有左陪集组成的集合，则 $H$ 还可在 $S$ 上以左平移作用：
  \[
    (h,xK)\longmapsto hxK.
  \]

  \item 设 $H\leq G$。可以定义 $H$ 在集合 $G$ 上的\textbf{共轭作用}
  \[
    (h,x)\longmapsto hxh^{-1}.
  \]
  由固定的 $h\in H$ 给出的映射 $x\mapsto hxh^{-1}$ 称为由 $h$ 带来的共轭，元素 $hxh^{-1}$ 称为 $x$ 的一个\textbf{共轭元}。

  若 $K\leq G$，则对任意 $h\in H$，都有
  \[
    hKh^{-1}\leq G,\qquad hKh^{-1}\cong K.
  \]
  因此 $H$ 也可以在 $G$ 的所有子群组成的集合上以共轭作用：
  \[
    (h,K)\longmapsto hKh^{-1}.
  \]
  子群 $hKh^{-1}$ 称为 $K$ 的一个\textbf{共轭子群}。
\end{enumerate}
\end{example}

\begin{purplenote}
在上述第 (3) 种作用中，为避免与第 (2) 种左平移的记号 $hx$ 混淆，通常直接写出 $hxh^{-1}$。
\end{purplenote}

% 整理：原稿 2.4.3(2) 的“x∈G”按作用的定义域修正为“x∈S”。
\begin{theorem}\label{thm:action-orbit-equivalence-stabilizer}
设群 $G$ 作用于集合 $S$，则：
\begin{enumerate}[label=(\arabic*)]
  \item 在 $S$ 上定义关系
  \[
    x\sim x'\quad\Longleftrightarrow\quad
    \text{存在 }g\in G\text{，使 }gx=x',
  \]
  则 $\sim$ 是等价关系。
  \item 对任意 $x\in S$，集合
  \[
    G_x=\{g\in G\mid gx=x\}
  \]
  是 $G$ 的一个子群。
\end{enumerate}
\end{theorem}

\begin{purplenote}
定理~\ref{thm:action-orbit-equivalence-stabilizer} 第 (1) 条中的等价类称为 $G$ 在 $S$ 上的\textbf{轨道}，$x\in S$ 所在的轨道记为
\[
  \bar x=\{gx\mid g\in G\}.
\]
$G_x$ 称为 $x$ 的\textbf{稳定子群}，也称\textbf{迷向群}。
\end{purplenote}

% 来源：微信图片_20261118134405_56_26.jpg，2.4.4--2.4.8 及紫笔说明、短证。
\begin{example}\label{ex:conjugacy-centralizer-normalizer}
若 $G$ 通过共轭作用于自身，则 $x\in G$ 的轨道
\[
  \bar x=\{gxg^{-1}\mid g\in G\}
\]
称为 $x$ 的\textbf{共轭类}。

若 $H\leq G$ 通过共轭作用于 $G$，则 $x$ 的稳定子群是
\[
  H_x=\{h\in H\mid hxh^{-1}=x\}
      =\{h\in H\mid hx=xh\}.
\]
它称为 $x$ 在 $H$ 中的\textbf{中心子}，记为 $C_H(x)$。当 $H=G$ 时，$C_G(x)$ 简称为 $x$ 的中心子。

若 $H$ 通过共轭作用于 $G$ 的所有子群组成的集合，则子群 $K\leq G$ 的稳定子群为
\[
  H_K=\{h\in H\mid hKh^{-1}=K\},
\]
称为 $K$ 在 $H$ 中的\textbf{正规子}，记为 $N_H(K)$。当 $H=G$ 时，$N_G(K)$ 简称为 $K$ 的正规子。
\end{example}

\begin{purplenote}
显然有
\[
  K\trianglelefteq N_G(K),\qquad
  \text{且有} K\trianglelefteq G\quad\Longleftrightarrow\quad N_G(K)=G.
\]
\end{purplenote}

\begin{theorem}\label{thm:orbit-stabilizer-cardinality}
设群 $G$ 作用于集合 $S$，则对每个 $x\in S$，其轨道的基数为
\[
  |\bar x|=[G:G_x].
\]
\end{theorem}

\begin{purpleproof}
对任意 $g,h\in G$，
\[
  gx=hx
  \quad\Longleftrightarrow\quad g^{-1}hx=x
  \quad\Longleftrightarrow\quad g^{-1}h\in G_x
  \quad\Longleftrightarrow\quad gG_x=hG_x.
\]
故 $gG_x\mapsto gx$ 给出了左陪集集合与轨道 $\bar x$ 之间的双射。
\end{purpleproof}

% 整理：类方程对全部共轭类求和，明确原稿“不同共轭类”是完整代表系。
\begin{corollary}\label{cor:conjugacy-class-equation}
设 $G$ 是有限群，$K\leq G$，则：
\begin{enumerate}[label=(\arabic*)]
  \item 对任意 $x\in G$，其共轭类中的元素个数为
  \[
    |\bar x|=[G:C_G(x)],\qquad [G:C_G(x)]\mid |G|.
  \]
  \item 若 $\bar x_1,\ldots,\bar x_n$ 是 $G$ 的全部不同共轭类，则
  \[
    |G|=\sum_{i=1}^{n}[G:C_G(x_i)].
  \]
  \item $G$ 中与 $K$ 共轭的子群的数目为
  \[
    [G:N_G(K)],\qquad [G:N_G(K)]\mid |G|.
  \]
\end{enumerate}
\end{corollary}

\begin{purplenote}
推论~\ref{cor:conjugacy-class-equation} 第 (2) 条的等式称为有限群 $G$ 的\textbf{类方程}。
\end{purplenote}

\begin{theorem}\label{thm:action-induced-permutation-homomorphism}
若群 $G$ 作用于集合 $S$，则这个作用诱导出一个群同态
\[
  \tau:G\longrightarrow A(S),
\]
其中 $A(S)$ 是 $S$ 上的所有置换组成的群。
\end{theorem}

% 整理：按原稿给出的 tau_g 定义，写明其逆映射及由作用公理得到的复合等式。
\begin{purpleproof}
对每个 $g\in G$，定义
\[
  \tau_g:S\longrightarrow S,\qquad \tau_g(x)=gx.
\]
由于 $\tau_{g^{-1}}$ 是 $\tau_g$ 的逆映射，$\tau_g\in A(S)$；又由作用公理，
\[
  \tau_{gh}(x)=(gh)x=g(hx)=(\tau_g\circ\tau_h)(x).
\]
故 $g\mapsto\tau_g$ 就是所需的群同态。
\end{purpleproof}

\begin{purplenote}
需注意这里并没有具体指定$G$在$S$上的作用映射，因此所诱导的群同态也不唯一。后续研究中，通常只需要知道存在这样的群同态即可。
\end{purplenote}

\begin{corollary}[Cayley]\label{cor:cayley-permutation-embedding}
对任意群 $G$，存在群单同态
\[
  G\longrightarrow A(G).
\]
% 根据 2026-11-19 的讨论恢复原稿“一些置换组成的群”的含义，明确为 A(G) 的子群。
因此任意群都同构于一个由某些置换组成的群，即 $A(G)$ 的一个子群。特别地，每个有限群 $G$ 都同构于 $S_n$ 的一个子群，其中 $n=|G|$。
\end{corollary}

\begin{purpleproof}
令 $G$ 在自身上以左平移作用，应用定理~\ref{thm:action-induced-permutation-homomorphism}。对任意 $x\in G$，由消去律有 $gx=x$ 当且仅当 $g=e$，故所诱导的群同态是单射。
\end{purpleproof}

% 来源：微信图片_20261118134408_57_26.jpg，2.4.9--2.4.12 及紫笔说明、证明。
\begin{purplenote}
对于群 $G$，$\operatorname{Aut}G$ 表示从 $G$ 到自身的所有群同构，即所有自同构组成的集合。以映射复合作为二元运算，$\operatorname{Aut}G$ 是一个群。
\end{purplenote}

\begin{corollary}\label{cor:inner-automorphisms-center-kernel}
设 $G$ 是群，则：
\begin{enumerate}[label=(\arabic*)]
  \item 对每个 $g\in G$，由 $g$ 诱导的共轭映射
  \[
    \tau_g:G\longrightarrow G,\qquad \tau_g(x)=gxg^{-1},
  \]
  是 $G$ 的自同构。
  \item 存在群同态
  \[
    \tau:G\longrightarrow\operatorname{Aut}G,\qquad g\longmapsto\tau_g,
  \]
  其核为
  \[
    C(G)=\{g\in G\mid gx=xg\text{ 对所有 }x\in G\text{ 成立}\}.
  \]
\end{enumerate}
\end{corollary}

\begin{purplenote}
推论~\ref{cor:inner-automorphisms-center-kernel} 第 (1) 条中的 $\tau_g$ 称为由 $g$ 诱导的\textbf{内同构}，正规子群 $C(G)=\ker\tau$ 称为 $G$ 的\textbf{中心}。

元素 $g\in C(G)$，当且仅当它的共轭类中只有 $g$ 一个元素。因此，当 $G$ 有限时，对 $x\in C(G)$ 有 $[G:C_G(x)]=1$，类方程可改写为
\[
  |G|=|C(G)|+\sum_{i=1}^{m}[G:C_G(x_i)],
\]
其中 $x_1,\ldots,x_m\in G\setminus C(G)$ 代表全部不同的非中心共轭类，各 $[G:C_G(x_i)]>1$。
\end{purplenote}

\begin{proposition}\label{prop:coset-action-kernel-contained-subgroup}
设 $G$ 是群，$H\leq G$。令 $G$ 在 $H$ 的所有左陪集组成的集合 $S$ 上以左平移作用，则所诱导的群同态
\[
  \tau:G\longrightarrow A(S)
\]
的核包含于 $H$。
\end{proposition}

\begin{purpleproof}
若 $g\in\ker\tau$，则 $g$ 固定左陪集 $H$，即 $gH=H$，故 $g\in H$。
\end{purpleproof}

\begin{corollary}\label{cor:core-free-subgroup-permutation-embedding}
设 $H\leq G$，$[G:H]=n$。若 $G$ 的非平凡正规子群均不包含于 $H$，则 $G$ 同构于 $S_n$ 的一个子群。
\end{corollary}

\begin{purpleproof}
直接应用命题~\ref{prop:coset-action-kernel-contained-subgroup}：所得同态的核是包含于 $H$ 的正规子群，由假设只能为 $\{e\}$。
\end{purpleproof}

\begin{corollary}\label{cor:smallest-prime-index-subgroup-normal}
设 $G$ 是有限群，$H\leq G$，且 $[G:H]=p$，其中 $p$ 是整除 $|G|$ 的最小质数，则 $H\trianglelefteq G$。
\end{corollary}

\begin{purpleproof}
令 $G$ 在 $H$ 的所有左陪集组成的集合 $S$ 上以左平移作用，记诱导的群同态为 $\varphi:G\to A(S)$，令 $K=\ker\varphi$。由命题~\ref{prop:coset-action-kernel-contained-subgroup}，$K\leq H$，只需证明 $K=H$。

因为 $|S|=p$，由\hyperref[cor:first-isomorphism]{第一同构定理}及\hyperref[cor:lagrange]{拉格朗日定理}，有
\[
  G/K\cong\operatorname{im}\varphi\leq A(S)\cong S_p,
  \qquad [G:K]\mid p!.
\]
又由\hyperref[thm:index-multiplication]{指数乘法公式}，
\[
  [G:K]=[G:H][H:K]=p[H:K],
  \qquad |G|=p[H:K]|K|.
\]
由 $p$ 的最小性，$[H:K]$ 的每个质因子都不小于 $p$；结合 $[G:K]\mid p!$，知这些质因子只能是 $p$，故 $[H:K]=p^a$，其中 $a\geq0$。但 $p!$ 中只含一个因子 $p$，而
\[
  p^{a+1}=[G:K]\mid p!,
\]
故 $a=0$，从而 $[H:K]=1$，即 $H=K\trianglelefteq G$。
\end{purpleproof}

% 来源：微信图片_20261118134410_58_26.jpg，2.4.13--2.4.15 及紫笔证明提示。
下面是本节的一些习题结论。

\begin{proposition}\label{prop:normalizer-centralizer-subgroup-properties}
设 $H\leq G$，则：
\begin{enumerate}[label=(\arabic*)]
  \item 若 $K\leq G$ 且 $H\trianglelefteq K$，则 $K\leq N_G(H)$。
  \item 定义
  \[
    C_G(H)=\{g\in G\mid gh=hg\text{ 对所有 }h\in H\text{ 成立}\},
  \]
  则 $C_G(H)\trianglelefteq N_G(H)$。
\end{enumerate}
\end{proposition}

% 字迹放大核对：原稿 2.4.14 是“阶为 p”的子群，不是“指数为 p”。
\begin{proposition}\label{prop:prime-order-subgroup-small-cofactor-normal}
若 $|G|=pn$，$p>n$ 且 $p$ 为质数，$H$ 是 $G$ 的一个阶为 $p$ 的子群，则 $H\trianglelefteq G$。
\end{proposition}

\begin{purplenote}
证明思路与推论~\ref{cor:smallest-prime-index-subgroup-normal} 的证明类似：仍考察陪集作用的核 $K$，证明 $H=K$ 即可。
\end{purplenote}

\begin{proposition}\label{prop:unique-normal-subgroup-prime-power-index}
若 $p,q$ 都是质数，$|G|=p^nq$ 且 $p>q$，则 $G$ 包含唯一一个指数为 $q$ 的正规子群。
\end{proposition}

\begin{purplenote}
% 整理引用：第 2.5 节已接入，为原稿“下一节”补充可点击的章节引用。
该结论要用到第~\ref{sec:sylow-theorems} 节的 Sylow 定理和推论~\ref{cor:smallest-prime-index-subgroup-normal}。
\end{purplenote}
